\begin{resumo}
Este trabalho tem como objetivo o desenvolvimento de um Sistema de Gerenciamento de Contratação de Estágios (SIGCE), 
focado na centralização e otimização do processo de assinatura e contratação de estágios na Faculdade de Ciências e Tecnologias em Engenharia (FCTE). 
O sistema proposto visa integrar e automatizar as etapas que envolvem a contratação e assinatura de estágios, melhorando a eficiência, acessibilidade e transparência das informações, 
além de reduzir a sobrecarga administrativa. 
Atualmente, o processo de contratação e assinatura de estágios na FCTE é descentralizado e realizado manualmente, 
o que gera problemas como a redundância de dados, inconsistência nas informações e dificuldades no acompanhamento das atividades. 
O sistema SIGCE propõe resolver essas questões ao centralizar as informações, melhorar a comunicação entre alunos, professores, 
orientadores e setores administrativos e garantir maior controle sobre os prazos e documentos obrigatórios, como termos de compromisso. 
O desenvolvimento do SIGCE adota uma metodologia ágil, utilizando o Scrum para garantir entregas incrementais e maior flexibilidade no processo. 
A arquitetura do sistema é baseada nos princípios de Arquitetura Limpa e Domain-Driven Design (DDD), com a implementação de microsserviços para garantir escalabilidade e flexibilidade. 
Além disso, o sistema integrará tecnologias como o Login Único do Governo Federal e a assinatura digital, garantindo maior segurança e conformidade com as normas legais. 
Ao final, espera-se que o SIGCE ofereça uma solução robusta para a gestão de contratação e assinatura de estágios na FCTE, promovendo uma experiência mais ágil e transparente para todos os envolvidos.

 \vspace{\onelineskip}
    
 \noindent
 \textbf{Palavras-chave}: assinatura digital. estágios. desenvolvimento. software.
\end{resumo}
